\emph{The quantitative headline summary is in \url{evaluation.md} §7,
Table 7. The narrative below is the systems-integration framing of what
those numbers mean.}

This work presented an effort to turn a lightweight open-source embedded
AI board into an open inference and autonomy platform for drones. The
central result is that an EdgeTPU-class accelerator can be paired
effectively with an MCU when the surrounding software stack is designed
to keep sensing, memory movement, scheduling, telemetry, and
programmability aligned. The goal is not to imitate a Linux-class
companion computer, but to show that a much lighter board can still
support meaningful onboard inference and autonomy when the system is
treated as a co-design problem rather than as a loose collection of
features.

The implementation extends the original Coral Micro board from which the
design started with a number of capabilities required for autonomous
drone research: a full MicroPython runtime, dual 5 Mpx cameras, onboard
IMU integration, a parallel camera-to-EdgeTPU perception pipeline,
CPU-side TensorFlow Lite Micro inference, lightweight tracking and
projection, runtime camera and model switching, interactive debugging
and deployment paths, obstacle-map reporting for PX4-style collision
prevention, dual-target control support spanning both Bitcraze Crazyflie
and PX4/MAVLink-style workflows, and an integrated set of classical ML
primitives that can be invoked directly from scripts. Taken together,
these additions define our MCU-EdgeTPU board as a more complete embedded
inference environment.

One of the clearest conclusions is that the hardest problems in compact
airborne AI systems often sit outside the neural model. Much of the
implementation effort went into ensuring that images could be acquired,
transferred, staged, quantized, and scheduled in real time without
breaking the low-power and low-mass profile of the platform. Host-side
systems engineering, especially memory placement, ownership discipline,
boot-time observability, and developer-facing tooling, proved just as
important as the accelerator.

The platform also contributes a broader methodological point. By
combining a task-oriented C++ runtime with a script-oriented MicroPython
layer, the system shows that it is possible to preserve real-time
discipline while still exposing an interface suitable for rapid
experimentation. This is particularly valuable in drone autonomy, where
sensing configuration, communication policies, flight behaviors, and
even the choice between accelerator-backed and MCU-only models are
rarely fixed early in development. The ability to move repeatedly
between low-level performance engineering and high-level scripting is
one of the main reasons the resulting system is useful as a research
artifact. The same interface now supports a wider set of embedded ML
experiments, including clustering, anomaly scoring, sequence modeling,
reinforcement learning, and lightweight SLAM, which further supports the
view of the runtime as an environment for embedded inference workflows
rather than a narrow detector wrapper.

Finally, the work suggests that lightweight open-source embedded AI
boards should not be viewed only as educational or prototyping tools.
With sufficient attention to integration quality, memory engineering,
runtime design, and flight-stack interfaces, they can serve as
technically relevant foundations for more demanding autonomous systems.
The resulting platform stands both as a concrete implementation and as
an argument for taking accessible open hardware seriously as a basis for
embedded autonomy research.
